\documentclass[../main.tex]{subfiles}
\begin{document}

\chapter{Discussion}
\label{ch:discussion}

\section{Interpreting the Detection Performance}

The Stage 1 evaluation results define the current functional boundaries of multi-agent LLM-based intrusion detection systems. Across all fourteen attack types, AMATAS achieves a mean F1 of 86.3\% (92.9\% excluding Infiltration), below the theoretical ceiling of fine-tuned discriminative models but substantially above both single-agent LLM approaches and, critically, the standalone Random Forest classifier. The RF baseline evaluation (Section~\ref{sec:rf-baseline}) reveals a striking pattern: the RF achieves perfect F1 on seven attack types with distinctive statistical signatures but 0\% on the remaining seven, yielding a mean F1 of only 57.1\%. AMATAS detects five of those seven RF-invisible categories, demonstrating that the multi-agent LLM pipeline's value lies not in matching ML accuracy on easy attacks but in providing detection capability where statistical classifiers fail entirely. The accuracy gap between AMATAS and fine-tuned models is the cost of generality and explainability.

Direct comparison with related work is instructive. \textcite{Houssel2024}, in the closest prior study, evaluated a standalone LLM-based NIDS and reported detection rates of approximately fifty per cent on benchmark flows, with significant false positive inflation. Houssel et al.\ characterise their evaluation as an ``initial exploration'' with limited sample sizes, so this figure should be interpreted as indicative rather than definitive; nonetheless, it establishes a baseline against which the multi-agent approach can be measured. AMATAS achieves roughly double that recall, a gain attributable to the shift from single-agent to multi-agent architecture, the adversarial validation provided by the Devil's Advocate, and the temporal context supplied by the cross-flow correlation agent. At the other extreme, \textcite{GutierrezGaleano2025} demonstrated F1 scores of 99.84 per cent using a fine-tuned T5 language model, though the authors acknowledge effectiveness limitations on attack types that are poorly represented in the dataset, which produced elevated false negative rates on minority classes. Their approach also produces no reasoning whatsoever; the fine-tuned model maps feature vectors to labels through learned weights that are as opaque as any traditional classifier. \textcite{Li2024} present IDS-Agent, a ReAct-style single-agent system reporting an F1 of 0.75 on the CIC-IoT dataset with iterative reasoning traces and tool-augmented analysis; AMATAS achieves F1=0.88 on a more comprehensive benchmark and extends this single-agent reasoning paradigm with multi-perspective specialist analysis, adversarial challenge, and consensus synthesis that distributes reasoning across six distinct analytical roles rather than a single iterative agent.

While AMATAS does not exceed the raw accuracy of fine-tuned discriminative models, it maintains an operationally viable detection rate while generating transparent, natural-language reasoning traces that traditional models do not provide. A significant constraint not addressed by the accuracy comparison is inference latency: \textcite{Kalafatidis2025} report that LLM processing consumed approximately 87\% of total execution time in their two-tier architecture (4.53 seconds per sample), and \textcite{Houssel2024} document a near-7000-fold slower inference time for LLMs compared to lightweight classifiers, concluding that real-time LLM-based NIDS deployment is not currently feasible. AMATAS inherits this limitation; the six-agent pipeline processes flows in seconds rather than microseconds, restricting its applicability to batch or near-real-time analysis rather than inline traffic inspection. The comparison with \textcite{Kalafatidis2025}, who independently propose a two-tier LLM-enhanced IDS architecture, further situates AMATAS within a converging research consensus, though AMATAS extends that pattern with a multi-agent debate structure and multi-perspective reasoning chains. Where \textcite{Kalafatidis2025} provide a single justification field alongside each threat assessment, AMATAS generates six distinct analytical narratives from specialist, adversarial, and synthesis perspectives, offering substantially richer explanatory depth for analyst review.

Within the Stage 1 results, the performance profile across attack types reveals important structural patterns. Brute-force and volumetric DoS categories (FTP-BruteForce (F1=100\%), SSH-Bruteforce (F1=99\%), DoS-Slowloris (F1=100\%), DoS-SlowHTTPTest (F1=100\%)) achieve the highest scores, reflecting the distinctive NetFlow signatures these attacks produce. \textcite{Leevy2020} note that perfect or near-perfect scores on the CICIDS2018 dataset are common across published studies and may indicate overfitting to dataset-specific patterns rather than genuine detection capability. The AMATAS results are partially insulated from this concern because the LLM agents are not trained on the dataset and reason from flow features at inference time; however, the possibility that the CICIDS2018 attack signatures are more distinctive than those in real-world traffic cannot be excluded, and validation on alternative datasets remains necessary. DDoS-HOIC (F1=72\%, recall=58\%) represents the hardest detectable category and is discussed at length in Section~\ref{sec:hoic}. Infiltration (F1=0\%, recall=0\%) is treated separately as a fundamental architectural limitation rather than a conventional performance shortfall.

\section{The Explainability Advantage}

This work demonstrates that while LLMs do not surpass traditional classifiers in raw accuracy, they offer competitive operational detection rates alongside detailed explanatory outputs that traditional models lack. Every AMATAS detection produces six distinct analytical outputs.

Consider flow 783 from the FTP-BruteForce batch as a representative example. The Protocol Agent classifies the flow as BENIGN at confidence 0.90, finding that the port assignment and protocol combination are individually consistent with legitimate FTP traffic. The Statistical Agent escalates to SUSPICIOUS at 0.75, flagging anomalous byte-per-packet ratios and inter-arrival timing. The Behavioural Agent returns SUSPICIOUS at 0.80, identifying a pattern consistent with credential stuffing but insufficient alone to confirm it. The Temporal Agent delivers MALICIOUS at confidence 0.95, having correlated the flow with 49 other connections from the same source IP in the analysis window; 49 connections to port 21 from a single source within a compressed time window is the decisive evidence that no single-flow analysis could provide. The Devil's Advocate acknowledges the temporal cluster but argues that the source IP could plausibly belong to a legitimate FTP client with an aggressive retry policy. The Orchestrator synthesises all five analyses into a SUSPICIOUS verdict at confidence 0.65 with consensus score 0.75, maps the finding to MITRE ATT\&CK technique T1110 (Brute Force), and produces a narrative explicitly reconciling the Protocol Agent's benign reading with the Temporal Agent's cluster evidence.

To make this contrast concrete, Table~\ref{tab:shap-vs-amatas} presents SHAP feature attributions from the Tier~1 Random Forest alongside AMATAS agent reasoning for five representative attack flows. The SHAP output for each flow is a ranked list of signed numeric values indicating which features pushed the classifier's prediction toward attack (positive) or benign (negative). The AMATAS output is a multi-perspective reasoning chain that contextualises those same features within protocol semantics, attack signatures, temporal patterns, and an explicit counter-argument. The Bot flow is particularly revealing: SHAP assigns exclusively negative (benign-leaning) values, consistent with the RF's P(attack) of 0.005, whilst AMATAS correctly identifies the flow as suspicious through temporal correlation with 49 other uniform flows from the same source IP. Similarly, for DDoS-HOIC, SHAP provides mixed signals with no clear attack indication, whilst AMATAS's Behavioural Agent identifies the specific pattern of short-duration HTTP sessions with asymmetric byte distributions characteristic of HTTP flood tools.

\begin{table}[H]
\centering
\caption{SHAP vs AMATAS: Five representative attack flows}
\label{tab:shap-vs-amatas}
\scriptsize
\begin{tabularx}{\textwidth}{@{}l l X@{}}
\toprule
\textbf{Flow} & \textbf{SHAP (RF)} & \textbf{AMATAS (most informative agent)} \\
\midrule
FTP-Brute\-Force (783) &
\texttt{L4\_DST\_PORT: +0.244} \newline \texttt{L4\_SRC\_PORT: +0.037} \newline \texttt{TCP\_FLAGS: +0.035} &
\textit{Temporal Agent} (MALICIOUS, 0.95): ``The flows from this source IP exhibit a pattern of repetitive, short-duration connections to the same destination IP and port (FTP service on port 21). Each flow is identical in duration, size, and TCP flags, suggesting automated behavior.'' \\
\addlinespace
SSH-Brute\-force (29) &
\texttt{L4\_DST\_PORT: +0.252} \newline \texttt{TCP\_FLAGS: +0.062} \newline \texttt{IN\_BYTES: +0.054} &
\textit{Temporal Agent} (MALICIOUS, 0.95): ``The temporal analysis of the flows from this source IP indicates a sustained and repetitive pattern targeting the same destination on port 22, indicative of a brute force SSH attack.'' \\
\addlinespace
DoS-Slow\-loris (763) &
\texttt{L4\_DST\_PORT: +0.134} \newline \texttt{TCP\_FLAGS: +0.089} \newline \texttt{L4\_SRC\_PORT: +0.072} &
\textit{Temporal Agent} (MALICIOUS, 0.95): ``The source IP shows a pattern of frequent, short-duration flows targeting the same destination IP and port (80). There is a notable burst activity with almost identical flows repeated multiple times.'' \\
\addlinespace
Bot (709) &
\texttt{L4\_DST\_PORT: \textbf{--0.208}} \newline \texttt{OUT\_BYTES: \textbf{--0.068}} \newline \texttt{IN\_BYTES: \textbf{--0.054}} &
\textit{Temporal Agent} (MALICIOUS, 0.95): ``The source IP has produced similar flows to the same destination and port with uniform packet sizes, flow durations, and TCP flags, indicating a potential automated tool.'' \\
\addlinespace
DDoS-HOIC (756) &
\texttt{OUT\_PKTS: \textbf{--0.126}} \newline \texttt{IN\_BYTES: \textbf{--0.069}} \newline \texttt{L4\_DST\_PORT: +0.061} &
\textit{Behavioural Agent} (MALICIOUS, 0.90): ``The flow involves HTTP (port 80) with very short duration sessions (16\,ms) and a higher number of outbound bytes compared to inbound bytes typical of data requests, consistent with HTTP-based flood tools.'' \\
\bottomrule
\end{tabularx}
\end{table}

This output is qualitatively different from what SHAP or LIME explanations provide for traditional classifiers. SHAP values tell an analyst which features influenced a model's output; they do not articulate what those features mean in the context of the specific flow, how they relate to known attack patterns, whether a reasonable benign interpretation exists, or how competing evidence was weighed. AMATAS produces all of these, in natural language accessible to analysts without machine learning expertise. This aligns with the chain-of-thought reasoning framework established by \textcite{Wei2022}, who demonstrate that eliciting intermediate reasoning steps substantially improves LLM performance on multi-step tasks, though they caution that there is no guarantee that generated reasoning paths are factually correct, that models can produce correct final answers through entirely incorrect reasoning, and that whether LLMs are genuinely ``reasoning'' remains an open question. This caveat applies directly to AMATAS, and the faithfulness audit (Section~\ref{sec:faithfulness-audit}) quantifies its extent: across 6{,}279 verifiable factual claims extracted from agent reasoning chains, 89.8\% were confirmed correct against actual flow data. Confabulation concentrated in TCP flag interpretation (22.4\% error rate) and protocol identification (19.4\%), where agents inferred expected values from domain knowledge rather than reading observed feature values. Numeric features (byte counts, packet counts, port numbers) were faithfully reported at greater than 98\% accuracy. This error profile is qualitatively distinct from open-ended hallucination: agents over-apply prior knowledge about protocol behaviour rather than fabricating unsupported content. The practical implication is that the reasoning traces are both inspectable and predominantly accurate: an analyst can read each agent's rationale, assess whether the cited evidence is compelling, and decide whether the conclusion follows, with confidence that the vast majority of factual claims are grounded in the actual flow data. \textcite{Hesford2024} underscore the urgency of this capability, demonstrating empirically that IDS performance degrades substantially between benchmark evaluation and real-world deployment, a gap that opaque classifiers give analysts no tools to diagnose. The practical implication is that AMATAS supports analyst decision-making rather than displacing it: an analyst reviewing a verdict can assess the reasoning, identify which agent analysis they find unconvincing, and apply domain expertise to confirm or override the system's conclusion.

A programmatic analysis of agent reasoning patterns across all 758 LLM-analysed flows confirms that the multi-agent architecture produces genuine analytical diversity rather than redundant commentary. Each agent exhibits a distinct feature focus profile: the Protocol Agent references port and protocol information on virtually every flow (99.5\% and 100\% respectively), while the Temporal Agent focuses on temporal patterns (85.2\%) and IP history (42.7\%), features that no other agent foregrounds. Pairwise agreement rates between specialist agents are notably low, with the protocol--temporal pair agreeing on only 24.1\% of flows, confirming that the agents analyse from genuinely different perspectives rather than producing correlated outputs. On the 671 flows where specialist verdicts were mixed, the Statistical Agent most frequently aligned with the final verdict (64.2\%), suggesting it serves as a de facto tiebreaker due to its focus on quantitative anomaly metrics. The Devil's Advocate exhibited systematically higher confidence on false negatives (mean 0.90) than on true positives (mean 0.80), providing further evidence that the fixed adversarial weight disproportionately suppresses correct detections on ambiguous flows --- consistent with the ablation finding that DA removal improves DDoS-HOIC recall from 58\% to 98\%.

\section{The False Positive Problem Resolved}

The zero per cent false positive rate across thirteen of fourteen attack types is perhaps the most practically significant result of this evaluation. Alert fatigue, the systematic desensitisation of security analysts to alerts generated by over-sensitive detection systems, is one of the most well-documented operational problems in network security. \textcite{Alahmadi2020} report that analysts in high-volume environments routinely ignore substantial fractions of generated alerts, with the consequence that genuine threats embedded within alert storms go uninvestigated. Alahmadi et al.\ further note that ``false positive'' is itself a broad metric in operational settings, distinguishing between true false alarms and ``benign triggers'' where the detection is technically correct but disregarded for business reasons. AMATAS's zero per cent FPR is measured against ground-truth labels rather than operational triage decisions, and real-world FPR perception would depend on the deployment context. However, \textcite{Alahmadi2020} caution that the term ``false positive'' is frequently conflated with ``benign trigger'' in operational contexts: many alerts that analysts classify as false positives are in fact true signature matches (e.g., detecting a vulnerable software version) that the organisation chooses to suppress for business reasons. Real-world alert fatigue is therefore driven partly by organisational context, not solely by model inaccuracy, and a 0\% FPR on a benchmark dataset does not guarantee elimination of alert fatigue in production deployment. Nonetheless, reducing technologically generated false positives remains a necessary condition for operational viability, even if it is not sufficient. At the CICIDS2018 dataset scale of twenty million flows, even a one per cent false positive rate would generate two hundred thousand spurious alerts, a volume that would render any system operationally useless.

The mechanism responsible for AMATAS's false positive suppression is the Devil's Advocate agent operating at a weight of thirty per cent in the Orchestrator's consensus calculation. This weight was calibrated through empirical ablation: at fifty per cent (Phase 3e), the agent became excessively aggressive, depressing recall below acceptable thresholds; at thirty per cent, it applies sufficient adversarial weighting to prevent the specialist agents from over-classifying ambiguous flows as malicious, without overwhelming genuinely malicious signals. This finding provides the first demonstration of the adversarial agent mechanism, described conceptually by \textcite{Yin2024} and the multi-agent debate literature \autocite{Du2023, Chan2023}, applied specifically to network intrusion detection with quantified performance consequences across a fourteen-category evaluation.

The ablation study (Section~\ref{sec:agent-ablation}) adds nuance to this picture: the DA's fixed weighting reduces recall on DDoS-HOIC from 98 per cent (no DA) to 58 per cent, indicating that on attacks with high feature overlap with benign traffic, the adversarial counter-argument carries disproportionate weight. The appropriate response is not to remove the DA but to make its weighting adaptive, a design change that would preserve its FPR control benefits on unambiguous attacks whilst reducing its suppressive effect on high-ambiguity categories.

\section{Limitations and Failure Modes}
\label{sec:limitations}

\subsection{The Infiltration Failure}

AMATAS failed to detect any of the fifty attack flows in the Infiltration category. This indicates an inherent limitation of flow-level analysis rather than a deficiency in model selection or prompt design. Infiltration attacks in the CICIDS2018 dataset are implemented via a simulated remote code execution chain that, at the individual flow level, generates traffic indistinguishable from normal HTTPS or SMB communication. There is no anomalous packet rate, no unusual port combination, no volume spike; the attacking flows conform in every measurable dimension to the statistical profile of the benign traffic in which they are embedded.

The solution this failure motivates is temporal clustering, the v3 architectural extension described in Section~\ref{sec:future}. The hypothesis is that providing the agent pipeline with ordered sequences of flows from the same source IP, rather than individual flows, would reveal the sequential access progression characterising infiltration attacks. Individual flows from an infiltrating host look benign; the ordered sequence of those flows (accessing a progression of internal services over a compressed time window) may not.

\subsection{DDoS-HOIC and Feature Overlap}
\label{sec:hoic}

DDoS-HOIC achieves recall of 58 per cent and F1 of 72 per cent, the lowest among successfully detected categories. The High Orbit Ion Cannon tool generates HTTP floods using randomised headers, user agents, and URL paths specifically to evade signature-based detection. The resulting flows share characteristics with high-volume legitimate web traffic: they use standard ports, conform to expected HTTP packet size distributions, and do not exhibit the periodicity that characterises simpler volumetric attacks. The Statistical Agent correctly identifies anomalous flow rates in 58 per cent of cases, but in the remaining 42 per cent the HOIC randomisation pushes the flow's statistical profile into the range occupied by legitimate web servers under genuine load.

\subsection{The Semantic Gap in Flow-Level Analysis}

Both the Infiltration failure and the HOIC difficulty point to a broader limitation identified by \textcite{Sommer2010}: the semantic gap between the statistical features of network traffic and their operational interpretation. Sommer and Paxson warn that network traffic features exhibit ``enormous variability'' and ``much more diversity than people intuitively expect,'' making it difficult to establish a stable notion of normality. They further caution that machine learning methods are better suited to finding similarities than to identifying activity that does not belong, a property they term the closed-world assumption. The AMATAS architecture partially addresses this gap by translating raw numerical features into natural language descriptions that the LLM can contextualise against its pre-trained knowledge of network protocols and attack patterns. However, this translation does not eliminate the underlying limitation: when malicious flows are statistically indistinguishable from benign traffic at the feature level, no amount of reasoning over those features can recover the missing signal. The Infiltration and HOIC results provide concrete evidence of this boundary.

\subsection{MCP Tool Utility and Anonymised Datasets}

The MCP comparison experiment (Section~\ref{sec:mcp-comparison}) evaluates whether external tool access improves single-agent detection. The results refute this hypothesis decisively: all four dataset-compatible tool configurations (D--G) performed worse than the no-tool baseline, despite receiving genuinely informative results from IANA, DShield, and CVE databases.

This result extends the tool-augmented LLM literature in an unexpected direction. \textcite{Yao2023} demonstrate that non-informative tool results derail reasoning; the AMATAS comparison demonstrates that \textit{informative} tool results can be equally detrimental when the information, though true, is irrelevant to the classification task. The qualitative analysis (Section~\ref{sec:mcp-qualitative}) reveals the mechanism: a normalising frame effect, consistent with the irrelevant context distraction documented by \textcite{Shi2023}, in which decoded feature values provide benign-interpretation templates that override the agent's anomaly detection. The agent anchors on tool-provided context at the expense of direct feature analysis, a form of reasoning displacement that \textcite{Qin2024} identify as a risk when tool outputs are more salient than the underlying data. The non-monotonic relationship between tool count and performance further suggests that additional tools compound the distraction rather than providing incremental value.

The substantial F1 gap between the best single-agent configuration and AMATAS (Section~\ref{sec:mcp-comparison}), combined with the finding that tool augmentation actively degrades performance, establishes that the multi-agent architecture's value derives from distributed specialist reasoning and adversarial validation rather than from external knowledge enrichment.

The agent ablation study (Section~\ref{sec:agent-ablation}) provides direct evidence that this multi-agent value is attack-type-dependent. The Temporal Agent emerges as the single most critical component, a result that aligns with the multi-agent debate framework of \textcite{Du2023}, who show that multiple model instances debating a question significantly outperform any single instance, with gains particularly pronounced on tasks requiring multi-step reasoning. The two-agent configuration's complete failure on SSH-Bruteforce confirms that multi-agent debate is not merely additive but constitutive of detection capability on attacks requiring temporal correlation. A related concern noted by both \citeauthor{Du2023} and \textcite{Chan2023} is that LLM performance can degrade as input context length increases, a risk structurally present in the Orchestrator agent, which receives all five preceding analyses.

The most significant ablation finding is the Devil's Advocate's non-monotonic contribution profile. On unambiguous attacks, DA removal produces no degradation; on high-ambiguity categories like DDoS-HOIC, the fixed adversarial weight suppresses correct but uncertain detections. The practical implication is confidence-gated adversarial weighting: the DA's influence should scale with specialist agreement rather than remaining fixed. This finding extends the multi-agent debate literature \autocite{Yin2024, Du2023, Chan2023} with the first empirical evidence that adversarial counter-argumentation can be counterproductive when classification difficulty is high. An important caveat from \textcite{Du2023} applies: multi-agent debates can converge on incorrect consensus, with agents confidently affirming a wrong answer. The Devil's Advocate partially mitigates this false consensus risk by structurally requiring a counter-argument regardless of specialist agreement, but it cannot detect cases where all specialists share the same incorrect premise.

The routing validation experiment (Section~\ref{sec:routing-validation}) establishes that the Tier~1 Random Forest contributes intelligent routing rather than random cost reduction, with implications beyond AMATAS: any two-tier system combining a cheap discriminator with an expensive analytical backend must demonstrate that the discriminator's selection is informed rather than arbitrary.

\subsection{Model Selection Constraints and Generalisability}

The exclusive reliance on GPT-4o means that reported performance figures may not generalise to other frontier LLMs. GPT-4o-mini exhibited a systematic false positive bias, classifying every flow as malicious regardless of content, which rendered it unusable for NIDS despite its cost advantage. Claude Sonnet 3.5 was rejected due to rate limits incompatible with 1000-flow batch processing and an estimated cost of approximately \$108 per batch. The multi-agent architecture is model-agnostic in principle, but the specific calibration of agent prompts, confidence thresholds, and Devil's Advocate weights was developed against GPT-4o behaviour. Alternative models may exhibit different reasoning patterns, hallucination profiles, and confidence tendencies requiring re-calibration.

The CICIDS2018 dataset, whilst the most comprehensive publicly available benchmark \autocite{Sharafaldin2018, Leevy2020}, is a simulation of network traffic generated in a controlled laboratory environment. The FLOW\_DIRECTION feature was unavailable in the CICIDS2018 NetFlow v3 encoding and was excluded from the Stage 1 feature set, which may limit discriminative power on directionally distinctive attacks. \textcite{Sarhan2021} note several limitations of the NetFlow-based representation, and \textcite{FengSakurai2025} contextualise dataset representativeness limitations within the broader challenge of constructing reliable NIDS evaluation benchmarks.

\section{Security Considerations for MCP-Enhanced Detection Systems}
\label{sec:mcp-security}

The MCP comparison experiments demonstrate that external tool access provides negligible benefit on anonymised datasets, but a production deployment of MCP-enhanced intrusion detection against live network traffic would face a qualitatively different threat model. The security properties of the MCP architecture itself (the protocol through which an LLM agent discovers, invokes, and trusts external tools) become critical concerns when the system operates on real traffic with genuine IP addresses, domain names, and payload content.

The most immediate risk is tool injection, where a compromised or malicious MCP server returns fabricated results designed to influence the LLM agent's classification. An attacker who controls an upstream threat intelligence feed could supply false reputation data, reporting a known malicious IP address as benign, or conversely, poisoning the reputation of legitimate infrastructure to trigger false positives at scale. The LLM agent, lacking independent means to verify the provenance or integrity of tool responses, would incorporate the fabricated data into its reasoning as though it were authoritative. This vulnerability is structural rather than incidental: the MCP specification delegates trust to the tool server, and no mechanism within the current protocol provides for cryptographic verification of response integrity or cross-validation against independent sources. \textcite{Maloyan2025} identify three fundamental protocol-level vulnerabilities in MCP-based architectures: the absence of capability attestation (servers cannot prove they are authorised to provide specific tools), bidirectional sampling risks (servers can invoke the LLM to execute unintended actions), and implicit trust propagation across the client-server boundary. Their proposed ATTESTMCP framework addresses these through capability attestation certificates, message authentication, and origin tagging of tool responses, providing a concrete defence-in-depth model for securing MCP deployments in security-critical contexts.

Permission escalation presents a related concern. An MCP server designed to provide read-only access to threat intelligence databases could, through crafted tool descriptions or parameter manipulation, induce the LLM agent to execute actions beyond its intended scope (writing to configuration files, modifying firewall rules, or exfiltrating sensitive flow data to external endpoints). The risk is amplified in intrusion detection contexts where the system necessarily processes sensitive network telemetry, and where an escalation from analysis to action could enable an attacker to manipulate the very defences the system is meant to provide.

Data exfiltration through MCP tool calls represents a third category of risk. Each tool invocation transmits a structured request to an external server, and in a NIDS context these requests may contain IP addresses, port numbers, flow statistics, and other telemetry that an adversary could aggregate to map the defended network's topology. Even if individual tool calls reveal limited information, the aggregate of thousands of calls across an analysis session could provide substantial intelligence to a patient adversary controlling or monitoring the tool server.

These considerations do not invalidate the MCP architecture for production deployment, but they establish that any future integration of external tools into the AMATAS pipeline must be accompanied by a security layer that the current MCP specification does not provide. The AMATAS design decision to embed domain knowledge in agent prompts rather than in external tool calls is, in retrospect, not merely a consequence of the anonymised dataset's incompatibility with threat intelligence lookups, but a security-conservative architectural choice that eliminates the tool-mediated attack surface entirely for the current deployment context.

\subsection{Multi-Agent Pipeline Security Risks}
\label{sec:pipeline-security}

The six-agent LLM pipeline itself, independent of MCP tool access, introduces security considerations that would apply to any production deployment.

\textbf{Adversarial flow crafting.} An attacker aware that an LLM-based NIDS analyses flow features could craft flows with feature values designed to trigger benign reasoning paths. The faithfulness audit (Section~\ref{sec:faithfulness-audit}) reveals a concrete vulnerability: agents confabulate TCP flag interpretations at a 22.4\% rate, inferring expected flags from protocol context rather than reading the actual bitmask. An adversary who understands this tendency could set TCP flags to values that the agents systematically misinterpret, exploiting the gap between what the agent believes and what the flow actually contains. More broadly, the LLM's reliance on pre-trained protocol knowledge means that flows deviating from textbook patterns in ways the model has not encountered may be systematically misclassified.

\textbf{Prompt injection via flow features.} Although the AMATAS agents receive flow features as structured numerical data rather than free text, the features are serialised into the prompt as key-value pairs. An attacker controlling source or destination ports could embed values that, when serialised, resemble natural language instructions (e.g., a port number that, combined with surrounding template text, creates a substring that alters the agent's interpretation). This risk is mitigated by the fixed prompt template and the numerical nature of NetFlow features, but it cannot be entirely excluded without formal input sanitisation.

\textbf{Cost-exhaustion attacks.} The two-tier architecture's economic viability depends on the Tier~1 Random Forest filtering approximately 95\% of flows. An attacker who crafts flows that consistently bypass the RF threshold (attack probability $\geq$ 0.15) without being genuinely malicious could force the expensive six-agent pipeline to process a disproportionate volume of benign traffic, exhausting the LLM budget without triggering any true detections. The \$5.00 per-batch cost limit provides a hard ceiling, but a sustained campaign could render the system economically non-viable.

\textbf{Consensus manipulation.} The Orchestrator synthesises five agent outputs into a verdict. An adversary who understands the consensus mechanism could craft flows that produce a specific split among specialist agents --- for example, flows that the Protocol and Statistical Agents classify as benign while the Behavioural and Temporal Agents flag as suspicious, forcing the Orchestrator into a low-confidence verdict that the Devil's Advocate can then suppress. The ablation study's finding that DDoS-HOIC produces exactly this pattern of moderate specialist confidence and DA over-suppression suggests that this attack surface is not hypothetical.

These risks are acknowledged rather than resolved in the current work. They apply to any LLM-based security system and are not specific to the AMATAS architecture, but they must be addressed before any production deployment.

\section{Threats to Validity}
\label{sec:threats-to-validity}

With respect to internal validity, each attack-type batch is constructed from a single random seed, meaning that performance figures reflect one specific sample of flows per category. A different seed might produce a batch with a different distribution of edge cases, potentially altering recall figures for categories with moderate sample-level variance. As \textcite{Cawley2010} emphasise, evaluation on a single partition may arbitrarily favour one configuration over another; multi-seed evaluations establishing confidence intervals are necessary to quantify this variance and are identified as a priority for future work.

With respect to construct validity, AMATAS produces three-class verdicts (BENIGN, SUSPICIOUS, and MALICIOUS), but the evaluation framework collapses SUSPICIOUS and MALICIOUS into a single detected category for binary metric calculation. This collapse is justified operationally (a SUSPICIOUS verdict triggers investigation just as a MALICIOUS verdict does), but it means that reported recall figures include flows classified at relatively low confidence, and a stricter MALICIOUS-only evaluation would produce lower recall figures.

A further construct-validity concern arises from known labelling errors in the CSE-CIC-IDS2018 dataset. \textcite{Liu2022} conducted a large-scale manual audit of the dataset's attack categories and documented pervasive mislabelling, particularly in the web-attack classes. Their audit reports that flows marked as SQL Injection frequently contain no SQL injection payload and consist instead of legitimate navigation traffic, that the Brute-Force category on 23 February 2018 is 41.7 per cent mislabelled with a substantial proportion of the corruption attached to UDP port 500 (IKE) background traffic, and that the Infiltration category exhibits label corruption severe enough that the original labelling logic cannot be reliably reconstructed. Manual inspection of the balanced evaluation batches used in the AMATAS versus single-agent comparison confirms this pattern: the SQL Injection batch contains twenty of twenty-five attack-labelled flows on UDP/500, the Brute-Force XSS batch contains twenty-one of twenty-five such flows, and the Brute-Force Web batch is a mix of Telnet, SSH, and IKE flows that do not match the expected HTTP signature. Because the ground truth on these three categories is itself unreliable, the large margins observed for AMATAS over the single-agent baseline on SQL Injection ($\Delta$MCC~$=+0.75$), Brute-Force XSS ($+0.71$), and Brute-Force Web ($+0.25$) should be interpreted as partial artefacts of AMATAS's greater propensity to flag anomalous-looking traffic, rather than as evidence of superior web-attack reasoning. The remaining AMATAS gains --- on DoS Slowloris ($+0.41$), DDoS HOIC ($+0.36$), Bot ($+0.33$), DoS Hulk ($+0.31$), DoS GoldenEye ($+0.18$), and DDoS LOIC-HTTP ($+0.15$) --- occur on batches whose labelled flows match their expected protocol and port signatures (TCP/80 for the HTTP-based attacks and TCP/8080 for the Ares command-and-control traffic characteristic of the CICIDS2018 Bot category), and therefore remain interpretable as genuine benefit from multi-agent deliberation over single-agent reasoning. Re-evaluation against the corrected dataset released by \textcite{Liu2022} is identified as an immediate priority for follow-up work.

The multi-agent architecture introduces a context length concern identified by both \textcite{Chan2023} and \textcite{Du2023}: as debate context grows, LLM performance can stagnate or degrade, with models tending to focus only on the most recent inputs. In the AMATAS pipeline, the Orchestrator receives the combined outputs of five preceding agents, which can exceed several thousand tokens. Whilst the current sequential (non-iterative) architecture limits context accumulation compared to the multi-round debates studied by Chan et al.\ and Du et al., flows with unusually long agent outputs may approach the threshold where context length degrades synthesis quality, a risk that would increase with any future extension to iterative multi-round debate.

The economic projection assumes that the Tier 1 Random Forest's one-hundred-per-cent recall on the development partition generalises to evaluation batches, and that cost per LLM call remains constant across flow types. In practice, the Temporal Agent's cost varies with the number of co-IP flows available for context, and flows with more complex feature patterns generate longer reasoning chains with higher token counts.

With respect to external validity, the CICIDS2018 dataset was generated in a controlled laboratory environment, and \textcite{Kenyon2020} and \textcite{Ring2019} caution that traffic generated in academic cyber-ranges may not represent the variability of real-world corporate and public networks. Performance figures obtained on this benchmark may not transfer directly to production traffic, where attack signatures, benign traffic profiles, and class distributions differ from the laboratory conditions. Additionally, the exclusive reliance on GPT-4o means that reported results reflect the behaviour of a single model family. \textcite{Chan2023} note that their multi-agent debate findings were validated only on homogeneous groups of LLMs from the same model family, and the transfer of these results to heterogeneous model configurations remains an open question relevant to any future multi-model deployment of AMATAS.

\onlystandalone{\printbibliography}
\end{document}
